\documentclass{amsbook}

\newtheorem{theorem}{Theorem}[chapter]
\newtheorem{lemma}[theorem]{Lemma}

\theoremstyle{definition}
\newtheorem{definition}[theorem]{Definition}
\newtheorem{example}[theorem]{Example}
\newtheorem{xca}[theorem]{Exercise}

\theoremstyle{remark}
\newtheorem{remark}[theorem]{Remark}

\numberwithin{section}{chapter}
\numberwithin{equation}{chapter}

\begin{document}

\frontmatter

\title{Minqi's Solutions to the Exercises of Walter Rudin (1976), Principles of Mathematical Analysis (3rd Edition)}

\author{Minqi Pan\footnote{Website: http://www.minqi-pan.com}}

\date{\today}

\maketitle

\setcounter{page}{4}

\tableofcontents

\mainmatter

\chapter{The Real and Complex Number Systems}

\section{Exercise 1: Mixed Operations between Rational and Irrational Numbers Produce Irrational Numbers}

If $r$ is rational ($r\ne 0$) and $x$ is irrational, prove that $r+x$ and $rx$ are irrational.

\begin{proof}
Denote $\alpha=r+x, \beta=rx$.

Suppose $\alpha\in\mathbb{Q}$. Then $x=\alpha-r$ is also in $\mathbb{Q}$, because $\mathbb{Q}$ is a field that is closed with respect to additive inverse and addition. This contradicts with the presumption that $x\notin\mathbb{Q}$.

Suppose $\beta\in\mathbb{Q}$. Then $x=\beta/r$ is also in $\mathbb{Q}$, because $r\ne 0$ and $\mathbb{Q}$ is a field that is closed with respect to multiplicative inverse and multiplication. This contradicts with the presumption that $x\notin\mathbb{Q}$.
\end{proof}

\section{Exercise 2: $\sqrt{12}\notin\mathbb{Q}$}

Prove that there is no rational number whose square is $12$.

\begin{lemma}\label{l1}
$\forall n\in \mathbb{Z}, \forall p\in \{2, 3\}$,
\[
n^2\equiv 0 \implies n\equiv 0 \pmod p.
\]
\end{lemma}

\begin{proof}
We check the squares of every element of $\mathbb{Z}/2\mathbb{Z}=\{0,1\}$, the ring of integers modulo $2$.
\[0^2 = 0\]
\[1^2 = 1\]

Thus $n^2=0$ implies $n=0$ in $\mathbb{Z}/2\mathbb{Z}$ by exhaustion.

Similarly we check the squares of every element in $\mathbb{Z}/3\mathbb{Z}=\{0,1,2\}$, the ring of integers modulo $3$.
\[0^2 = 0\]
\[1^2 = 1\]
\[2^2 = 1\]

Thus $n^2=0$ implies $n=0$ in $\mathbb{Z}/3\mathbb{Z}$ by exhaustion.
\end{proof}

\begin{lemma}\label{l2}
There is no rational number whose square is $3$.
\end{lemma}

\begin{proof}
We now show that the equation
\begin{equation}\label{e1}
p^2=3
\end{equation}
is not satisfied by any rational $p$. If there was such a $p$, we could write $p=\frac{m}{n}$ where $m\equiv 0 \pmod{3}$ and $n\equiv 0 \pmod{3}$ are not both true. Let us assume this is done. Then \ref{e1} implies
\begin{equation}\label{e2}
m^2=3n^2,
\end{equation}
This shows that $m^2\equiv 0\pmod{3}$ and we have $m\equiv 0 \pmod{3}$ by Lemma~\ref{l1}. So $m^2\equiv 0\pmod{9}$. It follows that the right side of \ref{e2} is divisible by 9, so that $n^2\equiv 0\pmod{3}$, which implies that $n\equiv 0\pmod{3}$ again by Lemma~\ref{l1}.

The assumption that \ref{e1} holds thus leads to the conclusion that both $m\equiv 0 \pmod{3}$ and $n\equiv 0 \pmod{3}$ are true, contrary to our choice of $m$ and $n$. Hence \ref{e1} is impossible for rational $p$.
\end{proof}

\begin{proof}[Main Proof]
Suppose that there existed $p\in\mathbb{Q}$ such that $p^2=12$. Then there exists $m, n\in\mathbb{Z}$ such that
\[
p=\frac{m}{n} \implies 12=\frac{m^2}{n^2} \implies m^2=2\times 6n^2 \implies m^2\equiv 0 \pmod{2}.
\]
So it follows from Lemma~\ref{l1} that
\[
m\equiv 0 \pmod{2}.
\]

Therefore there exists $k\in\mathbb{Z}$ such that $m=2k$, and
\[
4k^2=12n^2 \implies k^2=3n^2 \implies \left(\frac{k}{n}\right)^2=3.
\]

This contradicts with Lemma~\ref{l2}.

\end{proof}

\end{document}
